% ---
% id: reductions/rem:povmreg
% kind: rem
% title: "A general POVM on the register, and its price"
% summary: >-
%   Keeping the register intact in the target would need a projective
%   measurement of it; filing each branch as a block of the squashed system
%   removes that hypothesis. The price is bookkeeping only: a direct sum
%   instead of a tensor product, the branch weight still transported as an
%   operator.
% prominence: standard
% status: accepted
% uses: [thm:mixture]
% mentions: [def:reduction-instance, foundations/def:sector-map]
% assumes: []
% refs: []
% provenance: newly-derived
% lean:
% ---

\begin{remark}
\label{rem:povmreg}
Theorem~\ref{thm:mixture} assumes only that $\{N_j\}$ is a POVM. A
construction that keeps the register $D$ intact in the target, squashing
$\sum_jM_{v|j}\otimes N_j$ to $\sum_jF_{v|j}\otimes N_j$ on $\Qsys_0D$, does
need $\{N_j\}$ to be projective, because its proof turns on
$\mathcal{P}_i(X\otimes N_j)=\delta_{ij}X\otimes N_j$ for the projections
$\mathcal{P}_i=(\idon{}\otimes N_i)(\cdot)(\idon{}\otimes N_i)$, and that
identity fails for a general POVM. Filing each branch as a \emph{block} of
$\Qsys$ instead of a value of a retained register removes the hypothesis. The
price is bookkeeping and nothing else: the target is a direct sum rather than
a tensor product, which is the \tref[block list]{block-list} form of
\xref{foundations/def:sector-map} and the form every reduction instance emits
(Definition~\ref{def:reduction-instance}\,(I2)); and
Theorem~\ref{thm:mixture}\,(iii) shows the branch weight is still
transported, as the operator $\idon{\Bsys_0}\otimes N_j$ rather than as the
number it is not.
\end{remark}
